\documentclass[12pt]{article}
\usepackage[utf8]{inputenc}
\usepackage[T1]{fontenc}
\usepackage{hyperref}
\title{Tsaangi Language Resources for K-12 Curriculum Development}
\author{Compiled by Claude Sonnet 5 (claude-sonnet-5) for Andr\'as Kornai}
\date{}
\begin{document}
\maketitle

\section{Introduction}
Tsaangi (ISO 639-3 tsa; Glottocode tsaa1242) is a Bantu language spoken in Gabon and the Republic of Congo. Speaker figures differ: Wikipedia gives 22,000 (2000) and the ASJP wordlist page gives 31,610. It should not be confused with Tsaayi (Teke-Tsaayi, ISO tyi), which belongs to West Teke. Little material was found.

\section{Primary Educational Resources}
No primers, textbooks or school programmes were found.

\section{Linguistic and Reference Materials}
\begin{itemize}
\item Tsaangi language (Wikipedia, a stub): \url{https://en.wikipedia.org/wiki/Tsaangi_language}
\item ASJP wordlist for Tsaangi: \url{https://asjp.clld.org/languages/TSAANGI}
\item Glottolog entry: \url{https://glottolog.org/resource/languoid/id/tsaa1242.bigmap.html}
\end{itemize}
No dictionary or grammar was found.

\section{Religious and Community Resources}
Scripture Earth lists Tsaangi scripture resources, including portions on Bible.is and a New Testament on YouVersion:
\url{https://www.scriptureearth.org/00eng.php?iso=tsa}

\section{Digital and Audio Resources}
Global Recordings Network lists recordings for Teke-Tsaayi, a different language (tyi). They are given here only to prevent confusion and should not be treated as Tsaangi material:
\url{https://globalrecordings.net/en/language/tyi}

\section{Curriculum Development Context}
The documentation is thin. The main written resource is scripture, and no standardized orthography or school use was found. Any curriculum would need new materials developed with the community.
\end{document}